% Part: history
% Chapter: biographies
% Section: rozsa-peter

\documentclass[../../../include/open-logic-section]{subfiles}

\begin{document}

\olfileid{his}{bio}{pet} 

\olsection{રોઝા પેટેર}

\olphoto{peter-rozsa}{રોઝા પેટેર}
 
રોઝા પેટેરનો જન્મ 17 ફેબ્રુઆરી 1905ના
રોજ હંગેરીના બુડાપેસ્ટમાં રોઝા પોલિત્ઝર
નામે થયો હતો. પુનરાવર્તી વિધેયો પરના
તેમના કાર્ય માટે તેઓ સૌથી વધુ જાણીતા
છે; પુનરાવર્તન સિદ્ધાંતનું ક્ષેત્ર ઊભું
થવામાં આ કાર્ય પાયાનું હતું.

પેટેરનો ઉછેર મુશ્કેલ રાજકીય સમયમાં
થયો: તેઓ કિશોરાવસ્થામાં હતાં ત્યારે
પ્રથમ વિશ્વયુદ્ધ ચાલી રહ્યું હતું.
તેમ છતાં તેઓ બુડાપેસ્ટની સંપન્ન
મારિયા તેરેઝિયા કન્યાશાળામાં ભણી
શક્યાં અને 1922માં ત્યાંનો અભ્યાસ
પૂરો કર્યો. પછી તેમણે બુડાપેસ્ટની
પાઝમાન્ય પેટેર યુનિવર્સિટીમાં
(જેનું પાછળથી એટ્વોશ લોરાન્ડ
યુનિવર્સિટી નામ પડ્યું) અભ્યાસ
કર્યો. પિતાના આગ્રહથી શરૂઆતમાં
રસાયણશાસ્ત્ર ભણ્યાં, પરંતુ પછી
ગણિત તરફ વળ્યાં અને 1927માં
સ્નાતક થયાં. માધ્યમિક શાળામાં
ગણિત ભણાવવાની લાયકાત હોવા
છતાં, મહામંદીની વિશ્વ અર્થતંત્ર
પરની અસરને કારણે તે સમયની
આર્થિક સ્થિતિ અત્યંત ખરાબ હતી.
ત્યારે પેટેરે ખાનગી ટ્યુટર અને
ગણિતનાં ખાનગી શિક્ષક તરીકે
છૂટક કામ કર્યાં. આખરે તેઓ
ગણિતમાં ઉચ્ચ અભ્યાસ માટે
ફરી યુનિવર્સિટીમાં આવ્યાં.
શરૂઆતમાં તેમણે સંખ્યા-સિદ્ધાંત
પર કામ કરવાની યોજના બનાવી
હતી; પરંતુ પોતાનાં પરિણામો
અગાઉ જ સાબિત થઈ ચૂક્યાં
હોવાનું જાણતાં તેમણે ગણિત
લગભગ છોડી જ દીધું હોત.
તેમને ગેડલના અપૂર્ણતા
પ્રમેયો પર કામ કરવા પ્રોત્સાહિત
કરાયાં અને, પોતાને જાણ્યા
વિના, ગેડલનાં કેટલાંક પરિણામો
જુદી રીતોથી સાબિત કર્યાં.
તેનાથી તેમનો આત્મવિશ્વાસ
પાછો આવ્યો. ડેવિડ હિલ્બર્ટના
પાયાકીય કાર્યક્રમથી પ્રેરાઈ
તેમણે પુનરાવર્તન સિદ્ધાંત પર
પોતાના પહેલા લેખો લખ્યા.
1935માં તેમણે પીએચ.ડી. મેળવી
અને 1937માં \emph{Journal of
  Symbolic Logic}નાં સંપાદક બન્યાં.

પેટેરના પ્રારંભિક લેખોને
પુનરાવર્તી વિધેયોના સિદ્ધાંતના
પાયાના પ્રદાન તરીકે વ્યાપક
માન્યતા મળી છે.
\cite{Peter1935a}માં તેમણે
પુનરાવર્તનના જુદા પ્રકારો
વચ્ચેના સંબંધની તપાસ કરી.
\cite{Peter1935b}માં તેમણે
બતાવ્યું કે પુનરાવર્તી રીતે
વ્યાખ્યાયિત એક ચોક્કસ વિધેય
આદિમ પુનરાવર્તી નથી. આથી
વિલ્હેલ્મ એકરમાનનું અગાઉનું
પરિણામ વધુ સરળ બન્યું.
પેટેરે આપેલું આ સરળ વિધેય
આજે ઘણી વાર એકરમાન વિધેય,
અને ક્યારેક વધુ યોગ્ય રીતે
એકરમાન--પેટેર વિધેય કહેવાય
છે. પુનરાવર્તી વિધેયોના
સિદ્ધાંત પરનું પ્રથમ પુસ્તક
તેમણે લખ્યું \citep{Peter1951}.

તેમના કાર્યના મહત્ત્વ અને
પ્રભાવ છતાં પેટેરને 1945 સુધી
પૂર્ણ સમયનું અધ્યાપનપદ મળ્યું
ન હતું. બીજા વિશ્વયુદ્ધમાં
હંગેરી પર નાઝીઓના કબજા
દરમિયાન યહૂદીવિરોધી કાયદાઓને
કારણે તેમને ભણાવવાની
મંજૂરી ન હતી. 1944માં સરકારે
બુડાપેસ્ટમાં યહૂદીઓ માટે
અલગ ઘેટો બનાવ્યો; તેને
શહેરના બાકીના ભાગથી
અલગ કરાયો હતો અને તેની
રખેવાળી સશસ્ત્ર પહેરેદારો
કરતા હતા. 1945માં તેની
મુક્તિ થાય ત્યાં સુધી
પેટેરને ત્યાં રહેવાની
ફરજ પડી. ત્યાર પછી
તેમણે બુડાપેસ્ટની શિક્ષક
તાલીમ કૉલેજમાં અને
1955થી એટ્વોશ લોરાન્ડ
યુનિવર્સિટીમાં અધ્યાપન
કર્યું. ગણિતમાં અકાદમિક
ડૉક્ટર બનનાર તેઓ
હંગેરીની પ્રથમ સ્ત્રી
ગણિતશાસ્ત્રી હતાં અને
હંગેરીની વિજ્ઞાન અકાદમીમાં
ચૂંટાનાર પણ પ્રથમ સ્ત્રી હતાં.

પેટેર ગણિતના ઉત્સાહી
શિક્ષક તરીકે જાણીતા હતા.
માત્ર વ્યાખ્યાન આપવાને
બદલે તેઓ વિદ્યાર્થીઓ
સાથે ગાણિતિક પ્રશ્નોનું
સ્વરૂપ અને સૌંદર્ય શોધવાનું
પસંદ કરતા. તેથી વિદ્યાર્થીઓ
તેમને સ્નેહથી “રોઝા માસી”
કહેતા. 1977માં 71 વર્ષની
ઉંમરે તેમનું અવસાન થયું.

\begin{reading}
વધુ જીવનચરિત્રાત્મક વાંચન માટે
\citep{Oconnor2014} અને
\citep{Andrasfai1986} જુઓ.
\citet{Tamassy1994}એ પેટેરની
ટૂંકી મુલાકાત લીધી હતી.
ગણિત વિશે રસપ્રદ વાંચન
માટે પેટેરનું પુસ્તક
\emph{Playing With Infinity}
\citep{Peter2010} જુઓ.
\end{reading}

\end{document}
